\chapter{从开关到逻辑门}

上一章只得到一件事：输入信号可以控制一条电流路径是否导通。现在把这个能力用起来，让多个输入共同决定一个输出。先不要急着背 AND、OR、NOT 的符号，先看两个很小的硬件需求。

第一个需求是安全允许灯。小设备只有在“保护盖合上”和“急停未按下”两个条件同时成立时，允许灯才亮。第二个需求是故障灯。只要温度报警或电流报警任意一个出现，故障灯就要亮。第三个需求是反向控制：某个按钮按下时，输出反而要变低。这三个需求分别逼出 AND、OR、NOT。

\section{一、真值表先把行为列完整}

门电路的第一步不是画线，而是把输入组合列完整。两个输入 A、B 只有四种组合。电路必须对每一种组合给出明确输出，而且每一行都要满足上一章的要求：输出节点不能悬空，不能上下冲突。

先看安全允许灯。A 表示保护盖合上，B 表示急停未按下。允许灯亮的条件很严格：少一个条件都不行。

\begin{longtable}{p{0.12\linewidth}p{0.18\linewidth}p{0.16\linewidth}p{0.42\linewidth}}
\toprule
A & B & 允许灯 & 为什么 \\
\midrule
0 & 0 & 0 & 盖子没合上，急停也不是可运行状态 \\
0 & 1 & 0 & 急停条件满足，但盖子没合上 \\
1 & 0 & 0 & 盖子合上了，但急停条件不满足 \\
1 & 1 & 1 & 两个安全条件同时满足 \\
\bottomrule
\end{longtable}

这张表就是 AND 的行为。它不是抽象数学定义，而是“两个条件都满足才允许”的硬件约束。若某个实现让 A=1、B=0 时灯也亮，那就不是 AND 近似得不够好，而是安全行为错了。

再看故障灯。A 表示温度报警，B 表示电流报警。故障灯的规则反过来：只要有一个坏消息出现，就要亮。

\begin{longtable}{p{0.12\linewidth}p{0.18\linewidth}p{0.16\linewidth}p{0.42\linewidth}}
\toprule
A & B & 故障灯 & 为什么 \\
\midrule
0 & 0 & 0 & 没有任何故障证据 \\
0 & 1 & 1 & 电流报警已经足够触发故障 \\
1 & 0 & 1 & 温度报警已经足够触发故障 \\
1 & 1 & 1 & 两个故障同时存在，当然要报警 \\
\bottomrule
\end{longtable}

这张表就是 OR 的行为。AND 和 OR 的差别不在符号长什么样，而在需求不同：一个要求所有条件同时成立，一个要求任意证据成立。

NOT 只有一个输入，但同样要列完整。若 A 表示“禁止信号”，输出 Y 表示“允许动作”，那么 A=0 时 Y=1，A=1 时 Y=0。它解决的是反向解释：输入出现时，输出反而关闭。

\section{二、把真值表翻译成开关路径}

真值表说清楚“应该怎样”，开关路径说明“为什么会这样”。用受控开关看 AND，最自然的结构是串联。两个开关串在同一条路径上，只有 A 和 B 都让自己的开关闭合时，整条路径才通。任意一个开关断开，路径就断开。

把四种输入逐行走一遍，AND 就不需要死记：

\begin{longtable}{p{0.1\linewidth}p{0.1\linewidth}p{0.34\linewidth}p{0.34\linewidth}}
\toprule
A & B & 串联路径状态 & 输出为什么是这个结果 \\
\midrule
0 & 0 & 两个开关都断开 & 没有完整导通路径，输出不能被拉成 1 \\
0 & 1 & A 断开，B 闭合 & 串联路径仍被 A 切断 \\
1 & 0 & A 闭合，B 断开 & 串联路径仍被 B 切断 \\
1 & 1 & 两个开关都闭合 & 从源到输出的路径完整，输出成立 \\
\bottomrule
\end{longtable}

OR 更像并联。A 控制一条路径，B 控制另一条路径。只要任意一路闭合，输出就能被拉向目标电平。只有两路都断开时，输出才没有这类导通路径。

NOT 的核心是互斥驱动。输入为 0 时，上拉路径导通、下拉路径断开，输出被拉成 1；输入为 1 时，上拉路径断开、下拉路径导通，输出被拉成 0。这里要特别注意：不是把一个输入“改名”为相反值，而是用这个输入选择输出节点接向高电平路径还是低电平路径。

\begin{pdfdiagram}[x=1cm,y=1cm]
  \node[hw label,text=BookNavy] at (1.8,2.35) {AND：串联};
  \node[hw block,minimum width=1.0cm] (andA) at (0.9,1.35) {A};
  \node[hw block,minimum width=1.0cm] (andB) at (2.35,1.35) {B};
  \node[hw label] (andout) at (3.25,1.35) {通路};
  \draw[BookRule,line width=0.55pt] (0.0,1.35) -- (andA.west);
  \draw[BookRule,line width=0.55pt] (andA.east) -- (andB.west);
  \draw[BookRule,line width=0.55pt,-{Latex[length=2mm,width=1.4mm]}] (andB.east) -- (andout.west);
  \node[hw label,align=center] at (1.65,0.55) {A、B 都闭合\\整条路径才通};

  \node[hw label,text=BookNavy] at (6.0,2.35) {OR：并联};
  \node[hw block,minimum width=1.0cm] (orA) at (5.25,1.65) {A};
  \node[hw block,minimum width=1.0cm] (orB) at (5.25,0.9) {B};
  \coordinate (orleft) at (4.3,1.275);
  \coordinate (orright) at (6.35,1.275);
  \draw[BookRule,line width=0.55pt] (orleft) -- (4.65,1.65) -- (orA.west);
  \draw[BookRule,line width=0.55pt] (orleft) -- (4.65,0.9) -- (orB.west);
  \draw[BookRule,line width=0.55pt] (orA.east) -- (5.95,1.65) -- (orright);
  \draw[BookRule,line width=0.55pt] (orB.east) -- (5.95,0.9) -- (orright);
  \draw[BookRule,line width=0.55pt,-{Latex[length=2mm,width=1.4mm]}] (orright) -- (7.0,1.275);
  \node[hw label,align=center] at (5.65,0.2) {任意一路闭合\\输出路径就成立};

  \node[hw label,text=BookNavy] at (9.6,2.35) {NOT：互斥驱动};
  \node[hw block,minimum width=1.15cm] (hi) at (9.6,1.85) {上拉};
  \node[hw block,minimum width=1.15cm] (lo) at (9.6,0.8) {下拉};
  \node[hw state,minimum width=0.9cm] (noty) at (11.0,1.325) {Y};
  \draw[hw bus] (hi.east) -- (noty.west);
  \draw[hw bus] (lo.east) -- (noty.west);
  \node[hw control,minimum width=0.9cm] (nota) at (8.15,1.325) {A};
  \draw[hw bus] (nota.east) -- (hi.west);
  \draw[hw bus] (nota.east) -- (lo.west);
  \node[hw label,align=center] at (9.6,0.05) {A 只能让其中\\一条路径导通};
\end{pdfdiagram}
\diagramnote{串联、并联、互斥上拉/下拉分别给出 AND、OR、NOT 的开关直觉。}

这张图只表达直觉，还不是完整晶体管级电路。真实门要同时安排上拉网络和下拉网络，保证每一种输入组合下输出都被明确拉高或拉低。只画“能让输出为 1 的路径”还不够，还必须说明输出为 0 时由谁负责拉低。

\section{三、NAND 和 NOR 更接近基础积木}

如果只从逻辑式看，AND 和 OR 很自然；但在开关网络里，NAND 和 NOR 往往更像基础积木。原因是它们天然带反相输出，更容易形成互补的上拉和下拉结构：当下拉网络导通时输出被拉低；当下拉网络不导通时，上拉网络负责把输出拉高。

以 NAND 为例。NAND 的输出只有在 A=1 且 B=1 时为 0，其余三种情况为 1。

\begin{longtable}{p{0.1\linewidth}p{0.1\linewidth}p{0.22\linewidth}p{0.46\linewidth}}
\toprule
A & B & NAND 输出 & 开关直觉 \\
\midrule
0 & 0 & 1 & 下拉串联路径断开，上拉负责拉高 \\
0 & 1 & 1 & A 断开，下拉路径不完整 \\
1 & 0 & 1 & B 断开，下拉路径不完整 \\
1 & 1 & 0 & A、B 都闭合，下拉路径完整，输出被拉低 \\
\bottomrule
\end{longtable}

这就是为什么“输出那一端倒着看”很重要。问一个门是否可靠，不只问逻辑表达式对不对，还要问每一行输入下：谁拉高，谁拉低，会不会悬空，会不会上下同时强驱动。NAND 好用，是因为它能用清楚的下拉条件和对应上拉条件覆盖所有输入行。

只用 NAND 就能构造 NOT、AND、OR：

\begin{lstlisting}
NOT A = A NAND A
A AND B = (A NAND B) NAND (A NAND B)
A OR B = (A NAND A) NAND (B NAND B)
\end{lstlisting}

这几行说明，硬件不需要为每种逻辑关系重新发明一种器件。只要有一类稳定、可复制、驱动能力可控的基础门，就能组合出更复杂的逻辑网络。真实芯片里会有多种标准单元，例如不同驱动强度的 NAND、NOR、反相器、缓冲器、选择器，但“复杂逻辑来自小门组合”这条线没有变。

\section{四、XOR 表示两个输入是否不同}

XOR 的输出在两个输入不同的时候为 1，相同的时候为 0。它一开始看起来比 AND、OR 难，因为它不是简单的串联或并联；但它解决的是一个非常具体的问题：检测两个 bit 是否不同。

\begin{longtable}{p{0.16\linewidth}p{0.16\linewidth}p{0.24\linewidth}p{0.34\linewidth}}
\toprule
A & B & A XOR B & 解释 \\
\midrule
0 & 0 & 0 & 两个输入相同 \\
0 & 1 & 1 & 两个输入不同 \\
1 & 0 & 1 & 两个输入不同 \\
1 & 1 & 0 & 两个输入相同 \\
\bottomrule
\end{longtable}

可以把 XOR 拆成两种会输出 1 的情况：A=1 且 B=0，或者 A=0 且 B=1。写成逻辑表达式就是：

\begin{lstlisting}
A XOR B = (A AND NOT B) OR (NOT A AND B)
\end{lstlisting}

这条式子不要只当代数记忆。它对应两条证据路径：第一条证明“只有 A 为 1”，第二条证明“只有 B 为 1”。最后的 OR 合并这两条证据。若两个都为 0，没有证据成立；若两个都为 1，两条“只有一个为 1”的证据也都不成立。

XOR 后面会反复出现。两个 bit 相加时，和位就是“两个输入是否不同”；比较两个 bit 时，XOR 可以告诉我们它们是否不同；再接一个 NOT，就能得到“两个输入是否相同”。所以 XOR 虽然比 AND、OR 复杂，但它不是花哨门，而是算术和比较的基础小部件。

\section{五、门不是理想数学函数}

门电路在纸上像瞬间给出输出的函数，真实硬件却至少有三个限制。第一，输入变化后，输出需要传播延迟才稳定。第二，一个输出能驱动的后级输入数量有限，后级太多会让切换变慢，甚至让电平余量变差。第三，门的输出必须重新形成清楚的 0/1，不能只靠一个偏弱路径勉强改变节点电压。

\begin{longtable}{p{0.2\linewidth}p{0.38\linewidth}p{0.32\linewidth}}
\toprule
限制 & 含义 & 后续影响 \\
\midrule
传播延迟 & 输出不会在输入变化的同一瞬间稳定 & 组合逻辑层数越多，稳定越慢 \\
扇出 & 一个输出只能可靠驱动有限后级 & 大网络需要缓冲器和分层布线 \\
电平恢复 & 输出要重新形成清楚 0/1 & 不能只追求逻辑表达式正确 \\
\bottomrule
\end{longtable}

看一个小反例。若把十几个门的输入全接到同一个弱输出上，真值表仍然可以写得很漂亮，但物理电路可能切换很慢。输出从 0 变 1 时，要给所有后级输入电容充电；如果下一级在它还没进入可靠高电平区间时就读取，就可能看到旧值或不确定值。于是后面讲组合逻辑时，不能只数门的个数，还要看最长路径、负载和稳定时间。

本章把第一章的受控开关变成了门。AND 来自串联条件，OR 来自并联证据，NOT 来自互斥上拉/下拉，NAND/NOR 提供更适合复制的基础积木，XOR 提供“是否不同”的判断。下一章会把多个门接成组合逻辑网络。那时每个电路都要同时满足两件事：逻辑上真值表正确，物理上能在规定时间内稳定，并且能继续驱动后级。

\begin{classicnote}
本章要用两张表验收同一个门：真值表证明逻辑行为，路径表证明每一行为什么会被拉成对应电平。

\begin{itemize}
\item 纸上实验：任选一个二输入门，列出 4 行真值表，再为每一行写出“上拉网络导通/断开、下拉网络导通/断开”。
\item 机器证据：合法门在稳定输入下不能让输出悬空，也不能让上拉和下拉强路径同时导通。
\item 组合证据：用 NAND 构造 NOT、AND、OR 时，不要只写代数式，还要说明输出是否经过了电平恢复。
\item 检查题：一个 XOR 电路真值表正确，但比 AND 慢很多，这个慢来自功能复杂度、门级层数、扇出，还是三者都有可能？
\end{itemize}

经典教材的写法会反复提醒：逻辑代数不是电路本身。代数式说明目标行为，门网络说明实现路径，延迟和负载说明这条路径能不能在真实机器里按时工作。
\end{classicnote}
